%% 自動生成（phys/gen_phys.py）。直接編集しないこと
\makeatletter
\phdefvol{1}{力学}{動く・回る・ゆがむ・流れる}{「なぜ動く？」「なぜ浮かぶ？」に、力で答えよう}
\phdefface{1}{動く}
\phdefface{2}{回る}
\phdefface{3}{ゆがむ}
\phdefface{4}{流れる}
\def\phFrontLine{クレーンも、橋も、水道も。力学の4つの顔で読みとく}
\phdefpart{1}{第1部}{動く：質点と質点系}{（力-1〜力-4）}{1}
\phdefpart{2}{第2部}{回る：剛体と機構}{（力-5〜力-7）}{2}
\phdefpart{3}{第3部}{ゆがむ：材料の力学}{（力-8〜力-11）}{3}
\phdefpart{4}{第4部}{流れる：流体の力学}{（力-12〜力-16）}{4}
\def\ph@whyA{エレベータが上へ動き出すと、体が重く感じるのはなぜ？}\def\ph@whyrA{力-2\hspace{0.8em}\bkpageno{86}}
\def\ph@whyB{鉄でできた大きな船が、なぜ水に浮かぶ？}\def\ph@whyrB{力-12\hspace{0.8em}\bkpageno{644}}
\def\ph@whyC{蛇口から落ちる水は、なぜ下へいくほど細くなる？}\def\ph@whyrC{力-13\hspace{0.8em}\bkpageno{702}}
\def\ph@whyD{スケートの選手が腕をたたむと、速く回るのはなぜ？}\def\ph@whyrD{力-6\hspace{0.8em}\bkpageno{322}}
\def\ph@whyE{鉄橋は、なぜ三角形を組み合わせてつくる？}\def\ph@whyrE{力-5\hspace{0.8em}\bkpageno{260}}
\def\ph@whypage{906}
\def\phBackCatch{気になる「なぜ？」から、\\ページをめくってみよう}
\makeatother
